\chapter{適用例}\label{cha:Application}

本章では、第\ref{cha:Proposal}章で作成したプロンプトを、適用した結果を示す。
適用対象として、「授業出席管理システム」を選定した。
本事例を通じて、プロンプトへの入力内容、および生成する2段階の出力（章立て構造と詳細説明）の具体像を示す。

\section{プロンプトへの入力}

本プロンプトの「開発するシステム情報」セクションに、システム概要を以下のように入力した。

\begin{figure}[H]
\begin{tcolorbox}[colback=white, colframe=black]
\begin{verbatim}
システム名: 授業出席管理システム
概要: 学生が自身のスマートフォンを用いて授業の出席登録を行い、教員がリアルタイムで出席状況や過去の履歴を管理・集計するシステム。
\end{verbatim}
\end{tcolorbox}
\caption{適用例における開発するシステム情報}
\label{fig:app_input}
\end{figure}

\section{生成結果}

プロンプトを実行した結果、以下の2段階の出力を得た。

\subsection{章立て（構造）}

出力1は、章立て構造を示す。
画面仕様、操作フロー、機能仕様を見ると、ドメイン独自の章立て構造を生成していることが確認できる。


\begin{figure}[H]
\begin{tcolorbox}[
  colback=white,
  colframe=black,
  fonttitle=\bfseries
]
\begin{multicols}{2}
\footnotesize
\ttfamily
\fontsize{11}{12}\selectfont
\begin{verbatim}
1. 概要
    1.1. 目的
    1.2. システム概要
    1.3. 用語定義
    1.4. 適用範囲

2. システム全体像
    2.1. 利用者区分
    2.2. ユースケース一覧
    2.3. システム構成図

3. 動作環境・基本条件
    3.1. 動作環境
    3.2. ネットワーク条件
    3.3. 稼働条件
    3.4. 制約事項

4. 画面仕様
    4.1. 画面一覧
    4.2. 学生向け画面
        4.2.1. 出席登録画面
        4.2.2. 出席履歴確認画面
    4.3. 教員向け画面
        4.3.1. 出席状況確認画面
        4.3.2. 出席履歴管理画面

5. 操作フロー
    5.1. 学生操作フロー
    5.2. 教員操作フロー
    5.3. 例外操作フロー

6. 機能仕様
    6.1. 機能一覧
    6.2. 出席登録機能
    6.3. 出席状況確認機能
    6.4. 出席履歴管理・集計機能

7. データ仕様
    7.1. データ項目一覧
    7.2. データ型・制約
    7.3. 入出力データ定義

8. エラー・メッセージ仕様
    8.1. エラー一覧
    8.2. エラーメッセージ定義

9. 非機能要件
    9.1. パフォーマンス要件
    9.2. セキュリティ要件
    9.3. 可用性・信頼性要件

10. 運用・保守
    10.1. 通常運用
    10.2. 障害時運用
    10.3. データ保全

11. テスト観点整理
    11.1. テスト対象範囲
    11.2. テスト前提条件

12. 付帯事項
    12.1. 前提条件
    12.2. 制限事項
    12.3. 将来拡張の考慮
\end{verbatim}
\end{multicols}
\end{tcolorbox}
\caption{適用例における出力1（章立て構造）}
\label{fig:app_output1}
\end{figure}

\subsection{詳細説明付き章立て}

続いて、各章の内容を補足する詳細説明を出力した（図\ref{fig:app_output2}）。
全ての章、節、項に対して、記述指針を示している。
この詳細説明により、仕様書作成者は各章に何を記述すべきかを明確に理解できる。


% \begin{figure}[H]
\begin{tcolorbox}[
  colback=white,
  colframe=black,
  fonttitle=\bfseries,
  breakable
]
% \begin{multicols}{2}
\footnotesize
\ttfamily
\fontsize{11}{12}\selectfont
\begin{verbatim}
1. 概要：本システム仕様書の全体像と背景を説明する。
  1.1. 目的：本仕様書が開発者およびテスト担当者向けであることと利用目的を明確にする。
  1.2. システム概要：授業出席管理システムの機能概要を図と文章で説明する。
  1.3. 用語定義：本書で使用する専門用語や略語を初心者向けに定義する。
  1.4. 適用範囲：本仕様書が対象とする機能・画面・利用範囲を明示する。

2. システム全体像：システムを利用する立場と全体の関係性を整理する。
  2.1. 利用者区分：学生・教員などの利用者種別と役割を説明する。
  2.2. ユースケース一覧：利用者ごとの代表的な利用シナリオを整理する。
  2.3. システム構成図：システム全体の構成を簡易図で示し文章で補足する。

3. 動作環境・基本条件：システムが動作する前提条件を明確にする。
  3.1. 動作環境：対応OS、端末、ブラウザなどの動作環境を記載する。
  3.2. ネットワーク条件：インターネット接続や通信前提条件を説明する。
  3.3. 稼働条件：利用可能時間や同時利用の想定を記載する。
  3.4. 制約事項：システム上の制限や利用上の注意点を明示する。

4. 画面仕様：画面構成と画面ごとの内容を定義する。
  4.1. 画面一覧：本システムに存在する全画面を一覧化する。
  4.2. 学生向け画面：学生が利用する画面の仕様をまとめる。
    4.2.1. 出席登録画面：学生が授業出席を登録する画面のレイアウトと項目を示す。
    4.2.2. 出席履歴確認画面：学生が自身の出席履歴を確認する画面を説明する。
  4.3. 教員向け画面：教員が利用する画面の仕様をまとめる。
    4.3.1. 出席状況確認画面：教員が授業中の出席状況を確認する画面を説明する。
    4.3.2. 出席履歴管理画面：教員が過去の出席情報を管理・閲覧する画面を説明する。

5. 操作フロー：利用者の操作手順を時系列で示す。
  5.1. 学生操作フロー：学生が出席登録を行う一連の操作を説明する。
  5.2. 教員操作フロー：教員が出席状況を確認・管理する操作手順を説明する。
  5.3. 例外操作フロー：想定外操作やエラー時の利用者の流れを示す。

6. 機能仕様：システムが提供する外部機能の仕様を定義する。
  6.1. 機能一覧：本システムに含まれる全機能を一覧で示す。
  6.2. 出席登録機能：学生による出席登録の仕様と条件を説明する。
  6.3. 出席状況確認機能：教員がリアルタイムで出席状況を確認する機能を説明する。
  6.4. 出席履歴管理・集計機能：出席データの閲覧・集計に関する仕様を説明する。

7. データ仕様：外部から見えるデータ構造を定義する。
  7.1. データ項目一覧：扱う主要なデータ項目を一覧で示す。
  7.2. データ型・制約：各データの型、桁数、必須条件などを表で示す。
  7.3. 入出力データ定義：画面や機能で扱う入出力データを説明する。

8. エラー・メッセージ仕様：利用者に表示されるエラーを定義する。
  8.1. エラー一覧：想定されるエラー種別を整理する。
  8.2. エラーメッセージ定義：各エラー時に表示されるメッセージ内容を明示する。

9. 非機能要件：性能や安全性などの品質条件を定義する。
  9.1. パフォーマンス要件：応答時間や同時利用数などの要件を記載する。
  9.2. セキュリティ要件：認証、個人情報保護などの要件を説明する。
  9.3. 可用性・信頼性要件：障害耐性やデータ保護の考え方を示す。

10. 運用・保守：システム利用開始後の扱いを定義する。
  10.1. 通常運用：日常的な利用方法と管理内容を説明する。
  10.2. 障害時運用：障害発生時の対応方針を説明する。
  10.3. データ保全：データ保存やバックアップの考え方を示す。

11. テスト観点整理：テスト担当者が確認すべき観点を整理する。
  11.1. テスト対象範囲：本仕様書に基づくテスト範囲を明確にする。
  11.2. テスト前提条件：テスト実施時の前提条件や注意点を示す。

12. 付帯事項：仕様補足や将来への配慮事項をまとめる。
  12.1. 前提条件：本システム仕様の前提となる条件を整理する。
  12.2. 制限事項：現時点での仕様上の制限を明示する。
  12.3. 将来拡張の考慮：将来的な機能追加を想定した注意点を示す。
\end{verbatim}
% \end{multicols}
\end{tcolorbox}
\captionof{figure}{適用例における出力2（詳細説明付き章立て）}
\label{fig:app_output2}
% \end{figure}

\section{まとめ}

本適用例により、作成したプロンプトが章立てを生成できることが確認できた。
なお、次章の評価実験においては、実験の公平性と評価のしやすさを考慮し、評価対象を「出力1（章立て構造）」のみとしている。